\documentclass[aspectratio=169,11pt,t]{beamer}
\input{../../../../theme/ksa-cs-v2.tex}
% 원본 덱 (ksa-theme): slides/v2/ml1/week07/entropy/ -- 그림은 거기 figs/ 를 그대로 쓴다
\setbeamerfont{author}{size=\normalsize}
\setbeamerfont{institute}{size=\small}
\graphicspath{{../../../../v2/ml1/week07/entropy/figs/}{../../../../assets/}}
\makeatletter
\newcommand{\markpage}[1]{\protected@write\@auxout{}{\string\newlabel{#1}{{\the\c@framenumber}{}}}\hypertarget{#1}{}}
\newcommand{\pref}[1]{\hyperlink{#1}{\textcolor{ksagolddk}{(p.\ref{#1})}}}
\makeatother

\renewcommand{\ksacredit}{Machine Learning and Deep Learning \,\textperiodcentered\, 섀넌 엔트로피}
\title{섀넌 엔트로피 (Shannon Entropy)}
\subtitle{놀라움을 숫자로, 그리고 결정 트리의 정보 이득}
\author{Machine Learning and Deep Learning}
\institute{한국과학영재학교 (KSA)}
\date{}
\begin{document}
{\setbeamertemplate{footline}{}\begin{frame}[c,noframenumbering]\titlepage\end{frame}}

\section{질문}
% --- 오늘의 질문
\begin{frame}{오늘의 질문}
\vfill
\centering
{\Large 브라질이 한국을 이기면 \quad $\to$ \quad 그러려니}\\[0.8em]
{\Large 한국이 브라질을 이기면 \quad $\to$ \quad \kq{온 나라가 난리}}\\[1.4em]
{\large 같은 ``경기 결과 하나''인데 왜 무게가 다를까? \quad \textbf{놀라움을 숫자로 잴 수 있을까?}}
\vfill
\note{브라질이 한국을 이기면 사람들은 그러려니 합니다. 그런데 한국이 브라질을 이기면 온 나라가 난리가 납니다. 둘 다 경기 결과 하나를 알려 줄 뿐인데 왜 무게가 다를까요? 드문 일이 일어났기 때문입니다. 1948년 클로드 섀넌은 이 놀라움을 숫자로 재는 방법을 만들었고, 그게 오늘의 주제인 정보량과 엔트로피입니다. 지난 시간 결정 트리의 지니 대신 쓰는 섞임 측정이기도 하고, 2주차 교차 엔트로피의 뿌리이기도 합니다.}
\end{frame}

\section{놀라움}
% --- 놀라움
\begin{frame}{놀라움 (정보량): 드물수록 크다}
\markpage{surprise}
\vfill
\begin{columns}[c]
\begin{column}{0.5\textwidth}\centering\includegraphics[width=\linewidth]{en_surprise.png}\end{column}
\begin{column}{0.46\textwidth}
{\LARGE\[ I(x) = -\log_2 p(x) \]}
\renewcommand{\arraystretch}{1.3}
{\large\begin{tabular}{lr}
브라질 승 (90\%) & 0.15 비트\\
반반 (50\%) & 1 비트\\
한국 승 (10\%) & \kq{3.32 비트}\\
1\% 사건 & 6.64 비트
\end{tabular}}
\end{column}
\end{columns}
\vfill
\note{섀넌은 확률 p인 사건이 일어났을 때의 놀라움을 마이너스 로그 2 p로 정했습니다. 단위는 비트입니다. 90퍼센트 확률인 브라질 승리는 0.15비트, 거의 놀랍지 않습니다. 반반인 사건은 정확히 1비트입니다. 10퍼센트인 한국 승리는 3.32비트, 1퍼센트 사건은 6.64비트입니다. 확률이 작을수록 커지고, 확실한 일, 확률 1은 0비트입니다. 이 값을 정보량이라고도 부릅니다. 드문 결과일수록 우리가 몰랐던 것을 더 많이 알려 주니까요.}
\end{frame}

% --- 왜 로그
\begin{frame}{왜 하필 로그일까}
\vfill
\centering
{\large 독립인 두 사건이 함께 일어날 확률은 \textbf{곱} \quad 그런데 놀라움은 \textbf{더해져야} 자연스럽다}\vskip0.4em
{\LARGE\[ -\log_2 (p_1 p_2) = -\log_2 p_1 - \log_2 p_2 \]}\vskip0.4em
{\large 동전 한 번 앞면: 1 비트 \qquad 동전 세 번 모두 앞면 ($\frac18$): \kb{3 비트}}\\[0.8em]
{\large 1 비트 = 예 · 아니오 질문 하나의 답 \quad 8 가지 중 하나 = 질문 3 번 ($2^3 = 8$)}
\vfill
\note{왜 하필 로그일까요? 서로 독립인 두 사건이 함께 일어날 확률은 곱입니다. 그런데 놀라움은 더해지는 게 자연스럽습니다. 동전을 한 번 던져 앞면이 나온 놀라움에, 또 한 번 앞면이 나온 놀라움이 더해져야 하니까요. 곱을 합으로 바꾸는 함수가 로그입니다. 동전 한 번 앞면은 1비트, 세 번 모두 앞면은 확률 8분의 1이라 3비트입니다. 1비트는 예, 아니오 질문 하나의 답이 주는 정보입니다. 8가지 중 하나를 알아내려면 질문 3번이면 되니, 8분의 1 사건은 3비트입니다.}
\end{frame}

\section{엔트로피}
% --- 엔트로피
\begin{frame}{엔트로피 = 놀라움의 평균}
\markpage{entropy}
\vfill
\centering
{\LARGE\[ H = \sum_i p_i \cdot \big(-\log_2 p_i\big) = -\sum_i p_i \log_2 p_i \]}\vskip0.4em
{\large 한국 vs 브라질: \quad $0.9 \times 0.15 + 0.1 \times 3.32 = \kb{0.47}$ 비트}\\[0.8em]
{\large 경기를 보기 \textbf{전}, 결과가 평균적으로 얼마나 놀라울지 \quad = \quad 불확실함의 크기}
\vfill
\note{경기를 하기 전에, 결과가 평균적으로 얼마나 놀라울지 계산할 수 있습니다. 각 결과의 놀라움에 그 확률을 곱해 더하면 됩니다. 이것이 엔트로피입니다. 한국과 브라질 경기는 90퍼센트로 0.15비트, 10퍼센트로 3.32비트이니 평균 0.47비트입니다. 엔트로피는 결과를 보기 전의 불확실함의 크기라고 생각하면 됩니다. 결과가 뻔할수록 작고, 예측하기 어려울수록 큽니다.}
\end{frame}

% --- 곡선
\begin{frame}{팽팽할수록 엔트로피가 크다}
\vfill
\centering\includegraphics[width=0.6\linewidth]{en_binary.png}\par\vskip0.2em
{\large 조기축구 vs 브라질 (1\%): \textbf{0.08} \quad 한국 vs 브라질: 0.47 \quad 한일전 (반반): \kb{1.0, 최대}}\\[0.2em]
{\small 조기축구팀이 이기면 놀라움은 6.64 로 크지만, 1\% 라서 평균에는 거의 기여하지 않는다}
\vfill
\note{결과가 두 가지일 때 한쪽이 이길 확률에 따라 엔트로피를 그렸습니다. 조기축구팀과 브라질이면 조기축구팀이 이길 확률 1퍼센트, 엔트로피 0.08입니다. 이기면 놀라움은 6.64로 크지만 그럴 일이 거의 없으니 평균은 작습니다. 한국과 브라질은 0.47, 한일전처럼 반반이면 1로 최대입니다. 결과를 전혀 예측할 수 없을 때 엔트로피가 가장 큽니다. 한일전이 손에 땀을 쥐게 하는 이유입니다.}
\end{frame}

% --- 결과가 여러 개
\begin{frame}{결과가 여러 개면}
\vfill
\centering
\renewcommand{\arraystretch}{1.4}
{\Large\begin{tabular}{lcc}
경우 & 확률 & 엔트로피\\\hline
동전 & $\frac12, \frac12$ & 1\\
세 종 붓꽃 (50송이씩) & $\frac13, \frac13, \frac13$ & $\log_2 3 = 1.585$\\
주사위 & $\frac16 \times 6$ & $\log_2 6 = 2.585$\\
날씨 (맑음, 흐림, 비, 눈) & $\frac12, \frac14, \frac18, \frac18$ & \kb{1.75}
\end{tabular}}\vskip0.6em
{\large $k$ 가지가 모두 같은 확률이면 $\log_2 k$ 로 최대 \quad 고르지 않으면 더 작다}
\vfill
\note{결과가 여러 개여도 같은 식입니다. 세 종이 50송이씩인 붓꽃은 로그 2의 3, 1.585비트, 주사위는 2.585비트입니다. k가지가 모두 같은 확률이면 로그 2 k로 가장 큽니다. 날씨처럼 맑음이 절반, 흐림이 4분의 1, 비와 눈이 8분의 1씩이면, 결과는 4가지지만 엔트로피는 2가 아니라 1.75입니다. 맑을 거라고 꽤 예측할 수 있으니 불확실함이 줄어든 겁니다.}
\end{frame}

% --- 질문 수
\begin{frame}{엔트로피 = 평균 질문 수}
\vfill
\centering
{\large 날씨 맞히기 스무고개: \quad ``맑음?'' $\to$ ``흐림?'' $\to$ ``비?''}\vskip0.5em
\renewcommand{\arraystretch}{1.3}
{\Large\begin{tabular}{lccc}
날씨 & 확률 & 필요한 질문 & 확률 $\times$ 질문\\\hline
맑음 & $\frac12$ & 1 & 0.5\\
흐림 & $\frac14$ & 2 & 0.5\\
비 & $\frac18$ & 3 & 0.375\\
눈 & $\frac18$ & 3 & 0.375\\\hline
 & & 평균 & \kb{1.75}
\end{tabular}}\vskip0.4em
{\large 흔한 것부터 물으면 평균 질문 수 = 엔트로피 \quad (압축 · 부호화의 원리)}
\vfill
\note{엔트로피에는 멋진 해석이 있습니다. 날씨를 스무고개로 맞힌다고 해 봅시다. 흔한 것부터 묻습니다. 맑음이냐고 물어서 예면 1번에 끝나고, 아니면 흐림이냐고 물어 2번, 그것도 아니면 비냐고 물어 3번이면 비와 눈이 구별됩니다. 확률을 곱해 평균을 내면 1.75번, 엔트로피와 정확히 같습니다. 엔트로피는 결과를 알아내는 데 필요한 예, 아니오 질문 수의 평균 최솟값입니다. 파일 압축도 이 원리로, 자주 나오는 글자에 짧은 부호를 줍니다.}
\end{frame}

\section{결정 트리}
% --- 정보 이득
\begin{frame}{결정 트리에서: 정보 이득}
\markpage{ig}
\vfill
\centering
{\large \textbf{정보 이득} = 질문 전 엔트로피 $-$ 질문 후 엔트로피 (꽃 수로 가중 평균)}\vskip0.5em
{\Large\[ \text{꽃잎 길이} < 2.45: \quad 1.585 - \Big(\tfrac{50}{150}\cdot 0 + \tfrac{100}{150}\cdot 1\Big) = \kb{0.918} \]}
{\Large\[ \text{꽃잎 너비} < 1.75: \quad 1.585 - 0.899 = 0.686 \]}\vskip0.3em
{\large 지니로 고를 때와 \textbf{같은 질문}을 고른다 \quad 다음 질문 (너비 1.75): 정보 이득 0.69}
\vfill
\note{지난 시간 결정 트리에서 지니 대신 엔트로피를 쓸 수 있습니다. 붓꽃 세 종이 같은 비율이니 처음 엔트로피는 1.585비트입니다. 꽃잎 길이 2.45로 나누면 왼쪽 setosa 50송이는 0비트, 오른쪽은 두 종 반반이라 1비트, 가중 평균은 0.667비트입니다. 1.585에서 0.667을 빼면 0.918비트만큼 불확실함이 줄었습니다. 이 줄어든 양을 정보 이득이라고 부릅니다. 꽃잎 너비 1.75로 나누면 정보 이득은 0.686으로 더 작습니다. 그래서 지니와 똑같이 꽃잎 길이 질문을 고릅니다. 이 질문이 꽃의 종에 대해 가장 많은 정보를 준다는 뜻입니다.}
\end{frame}

% --- 지니 vs 엔트로피
\begin{frame}{지니와 엔트로피: 거의 같은 모양}
\vfill
\centering\includegraphics[width=0.56\linewidth]{en_gini.png}\par\vskip0.2em
{\large 둘 다 반반일 때 최대, 한 종만 있으면 0 \quad 지니 $\times 2$ 와 엔트로피는 거의 겹친다}\\[0.2em]
{\small scikit-learn \texttt{criterion="gini"} (기본) \;/\; \texttt{"entropy"} \quad 결과 트리는 대개 비슷하다}
\vfill
\note{두 종일 때 지니와 엔트로피를 같이 그려 봤습니다. 둘 다 반반일 때 가장 크고 한 종만 있으면 0입니다. 지니에 2를 곱하면 엔트로피 곡선과 거의 겹칩니다. 그래서 어느 것을 써도 결정 트리가 고르는 질문은 대개 같습니다. scikit-learn은 기본으로 지니를 쓰는데, 로그 계산이 없어서 조금 빠르기 때문입니다. criterion을 entropy로 바꿀 수 있습니다.}
\end{frame}

% --- 이어지는 곳
\begin{frame}{엔트로피는 어디로 이어지나}
\vfill
\centering
\renewcommand{\arraystretch}{1.6}
{\large\begin{tabular}{ll}
\textbf{교차 엔트로피} (2주차) & 모델 $q$ 의 놀라움으로 잰 평균 $-\sum p \log q$\\
\textbf{KL 발산} (15주차) & 교차 엔트로피 $-$ 엔트로피 = 모델이 틀려서 \textbf{더} 놀란 양\\
\textbf{언어 모델} (17주차) & 다음 글자 · 단어의 놀라움을 줄이도록 학습
\end{tabular}}\vskip0.8em
{\small 이름의 유래: AI 모델 ``Claude'' 는 섀넌의 이름 (Claude Shannon) 에서 따왔다고 알려져 있다}
\vfill
\note{엔트로피는 이 수업 곳곳으로 이어집니다. 2주차에 배운 교차 엔트로피는 진짜 분포의 사건을 모델의 확률로 잰 놀라움의 평균입니다. 모델이 정답에 낮은 확률을 주면 놀라움이 커서 벌점이 큽니다. 교차 엔트로피에서 엔트로피를 빼면 15주차의 KL 발산, 모델이 틀려서 더 놀란 양이 됩니다. 언어 모델은 다음 글자나 단어의 놀라움을 줄이도록 학습합니다. 참고로 AI 모델 클로드의 이름도 섀넌에게서 따왔다고 알려져 있습니다.}
\end{frame}

\section{연습}
% --- 문제풀이
\begin{frame}{문제풀이}
\vfill
\begin{enumerate}
\item 확률 $\frac14$ 인 사건의 놀라움은? 확률 $\frac{1}{1024}$ 이면?\vskip0.8em
\item 결과가 A 75\%, B 25\% 인 경기의 엔트로피는? ($\log_2 3 \approx 1.585$)\vskip0.8em
\item 사과 4 · 귤 4 인 칸을 (사과 4) / (귤 4) 로 나눴다. 정보 이득은?\vskip0.8em
\item 8가지 결과가 모두 같은 확률이면 엔트로피는? 이 결과를 알아내는 데 예 · 아니오 질문이 몇 번 필요한가?
\end{enumerate}
\vfill
\note{1번 답. 4분의 1은 2비트, 1024분의 1은 10비트입니다. 2번 답. 0.75 곱하기 마이너스 로그 2의 0.75, 이는 2 빼기 로그 2의 3, 약 0.415비트이고, 여기에 0.25 곱하기 2비트를 더하면 0.311 더하기 0.5로 약 0.811비트입니다. 3번 답. 처음 반반이라 1비트, 나눈 뒤 양쪽 모두 0비트이니 정보 이득 1비트입니다. 4번 답. 3비트, 질문 3번입니다.}
\end{frame}

% --- 실습
\begin{frame}{실습}
\vfill
\centering\begin{tabular}{rl}
\large 1 & 놀라움 · 엔트로피 함수 직접 구현 $\to$ 0.47, 1.75 재현\\[0.5em]
\large 2 & 정보 이득으로 결정 트리 질문 고르기, \texttt{criterion="entropy"} 와 비교\\[0.5em]
\large 3 & 실제 글: 셜록 홈스 소설의 글자 엔트로피 $\to$ 왜 압축이 되는가
\end{tabular}\vskip1em
The Adventures of Sherlock Holmes (Project Gutenberg) \\[0.4em]\href{https://colab.research.google.com/github/SuminHan/book-ml/blob/main/slides/v2/ml1/week07/entropy/lab/entropy.ipynb}{\textcolor{ksagolddk}{▶ Colab에서 실습 노트북 열기}}
\vfill
\note{실습은 세 부분입니다. 첫째, 놀라움과 엔트로피 함수를 직접 만들어 오늘의 0.47과 1.75를 재현합니다. 둘째, 정보 이득으로 붓꽃의 첫 질문을 고르고, scikit-learn에서 criterion을 entropy로 바꾼 트리와 비교합니다. 셋째, 셜록 홈스 소설 전체의 글자 빈도로 엔트로피를 재서, 영어 글자 하나에 평균 몇 비트가 필요한지, 그리고 왜 텍스트 파일이 잘 압축되는지 확인합니다.}
\end{frame}

\section{정리}
% --- 정리
\begin{frame}{정리}
\vfill
\centering\large\begin{tabular}{l}
\textbf{놀라움} $-\log_2 p$ \quad 드물수록 크다, 독립 사건이면 더해진다\\[0.5em]
\textbf{엔트로피} $H = -\sum p \log_2 p$ \quad 놀라움의 평균 = 불확실함 = 평균 질문 수\\[0.5em]
팽팽할수록 최대 (두 결과면 반반에서 1 비트, $k$ 가지면 $\log_2 k$)\\[0.5em]
\textbf{정보 이득} = 질문 전 $-$ 질문 후 엔트로피 \quad 결정 트리가 질문을 고르는 기준
\end{tabular}
\vfill
\note{정리합니다. 놀라움, 정보량은 마이너스 로그 확률로, 드물수록 크고 독립 사건이면 더해집니다. 엔트로피는 놀라움의 평균이고, 불확실함의 크기이자 결과를 알아내는 데 필요한 평균 질문 수입니다. 결과가 팽팽할수록 커서 두 결과면 반반일 때 1비트, k가지면 로그 2 k가 최대입니다. 질문 전후의 엔트로피 차이가 정보 이득이고, 결정 트리는 정보 이득이 가장 큰 질문을 고릅니다.}
\end{frame}
\end{document}
